\documentclass{amsart}
% For dealing with references we use the comment environment
\usepackage{verbatim}
\newenvironment{reference}{\comment}{\endcomment}
%\newenvironment{reference}{}{}
\newenvironment{slogan}{\comment}{\endcomment}
\newenvironment{history}{\comment}{\endcomment}

% For commutative diagrams we use Xy-pic
\usepackage[all]{xy}

% We use 2cell for 2-commutative diagrams.
\xyoption{2cell}
\UseAllTwocells

% We use multicol for the list of chapters between chapters
\usepackage{multicol}

% This is generally recommended for better output
\usepackage{lmodern}
\usepackage[T1]{fontenc}

% For cross-file-references

% Package for hypertext links:
\usepackage{hyperref}

% For any local file, say "hello.tex" you want to link to please

% Theorem environments.
%
\theoremstyle{plain}
\newtheorem{theorem}[subsection]{Theorem}
\newtheorem{proposition}[subsection]{Proposition}
\newtheorem{lemma}[subsection]{Lemma}

\theoremstyle{definition}
\newtheorem{definition}[subsection]{Definition}
\newtheorem{example}[subsection]{Example}
\newtheorem{exercise}[subsection]{Exercise}
\newtheorem{situation}[subsection]{Situation}

\theoremstyle{remark}
\newtheorem{remark}[subsection]{Remark}
\newtheorem{remarks}[subsection]{Remarks}

\numberwithin{equation}{subsection}

% Macros
%
\def\lim{\mathop{\mathrm{lim}}\nolimits}
\def\colim{\mathop{\mathrm{colim}}\nolimits}
\def\Spec{\mathop{\mathrm{Spec}}}
\def\Hom{\mathop{\mathrm{Hom}}\nolimits}
\def\Ext{\mathop{\mathrm{Ext}}\nolimits}
\def\SheafHom{\mathop{\mathcal{H}\!\mathit{om}}\nolimits}
\def\SheafExt{\mathop{\mathcal{E}\!\mathit{xt}}\nolimits}
\def\Sch{\mathit{Sch}}
\def\Mor{\mathop{\mathrm{Mor}}\nolimits}
\def\Ob{\mathop{\mathrm{Ob}}\nolimits}
\def\Sh{\mathop{\mathit{Sh}}\nolimits}
\def\NL{\mathop{N\!L}\nolimits}
\def\CH{\mathop{\mathrm{CH}}\nolimits}
\def\proetale{{pro\text{-}\acute{e}tale}}
\def\etale{{\acute{e}tale}}
\def\QCoh{\mathit{QCoh}}
\def\Ker{\mathop{\mathrm{Ker}}}
\def\Im{\mathop{\mathrm{Im}}}
\def\Coker{\mathop{\mathrm{Coker}}}
\def\Coim{\mathop{\mathrm{Coim}}}

% Boxtimes
%
\DeclareMathSymbol{\boxtimes}{\mathbin}{AMSa}{"02}

%
% Macros for moduli stacks/spaces
%
\def\QCohstack{\mathcal{QC}\!\mathit{oh}}
\def\Cohstack{\mathcal{C}\!\mathit{oh}}
\def\Spacesstack{\mathcal{S}\!\mathit{paces}}
\def\Quotfunctor{\mathrm{Quot}}
\def\Hilbfunctor{\mathrm{Hilb}}
\def\Curvesstack{\mathcal{C}\!\mathit{urves}}
\def\Polarizedstack{\mathcal{P}\!\mathit{olarized}}
\def\Complexesstack{\mathcal{C}\!\mathit{omplexes}}
% \Pic is the operator that assigns to X its picard group, usage \Pic(X)
% \Picardstack_{X/B} denotes the Picard stack of X over B
% \Picardfunctor_{X/B} denotes the Picard functor of X over B
\def\Pic{\mathop{\mathrm{Pic}}\nolimits}
\def\Picardstack{\mathcal{P}\!\mathit{ic}}
\def\Picardfunctor{\mathrm{Pic}}
\def\Deformationcategory{\mathcal{D}\!\mathit{ef}}

\setlength{\emergencystretch}{3em}
\begin{document}
\title[Stacks Project, Tag 0E4L]{255 --- Perfection, pushforward vanishing and self-duality of proper flat relative dualizing complexes}
\maketitle
\noindent Copyright (C) 2005--2025 Johan de Jong. Stacks Project authors. Source: \href{https://stacks.math.columbia.edu/tag/0E4L}{Tag 0E4L}.

\noindent Editorial difficulty note (not author text). Difficulty H2: Relative perfection is detected by tensoring against every perfect object and dualizing its pushforward. The proof then establishes pushforward amplitude and reconstructs the self-Hom unit through functorial test-object identities and generation by perfect complexes. These are multiple substantive properties of an adjoint-defined complex..

\noindent Editorial provenance note (not author text). History: excerpt from revision \texttt{a0444\allowbreak 6e57e\allowbreak c1fbc\allowbreak 252a8\allowbreak 71afc\allowbreak ec775\allowbreak 2fb28\allowbreak 07b14}, duality.tex, lines 2810--2939. Reference presentation was adapted; original-author context and core support blocks were retained or added. The mathematical wording is unchanged. Licensed under GNU FDL 1.2 or later; no invariant sections or cover texts. See ../licenses/COPYING. Context blocks reproduce author definitions and supporting results; they are not additional counted problems.

\begin{remark}
\label{remark-relative-dualizing-complex}
Let $Y$ be a quasi-compact and quasi-separated scheme.
Let $f : X \to Y$ be a proper, flat morphism of finite presentation.
Let $a$ be the adjoint of Lemma \href{https://stacks.math.columbia.edu/tag/0A9E}{lemma twisted inverse image (Tag 0A9E)} for $f$.
In this situation, $\omega_{X/Y}^\bullet = a(\mathcal{O}_Y)$
is sometimes called the {\it relative dualizing complex}. By
Lemma \href{https://stacks.math.columbia.edu/tag/0E4K}{lemma compare with pullback flat proper (Tag 0E4K)}
there is a functorial isomorphism
$a(K) = Lf^*K \otimes_{\mathcal{O}_X}^\mathbf{L} \omega_{X/Y}^\bullet$
for $K \in D_\QCoh(\mathcal{O}_Y)$. Moreover, the trace map
$$
\text{Tr}_{f, \mathcal{O}_Y} : Rf_*\omega_{X/Y}^\bullet \to \mathcal{O}_Y
$$
of Section \href{https://stacks.math.columbia.edu/tag/0AWG}{section trace (Tag 0AWG)} induces the trace map for all $K$
in $D_\QCoh(\mathcal{O}_Y)$. More precisely the diagram
$$
\xymatrix{
Rf_*a(K) \ar[rrr]_{\text{Tr}_{f, K}} \ar@{=}[d] & & &
K \ar@{=}[d] \\
Rf_*(Lf^*K \otimes_{\mathcal{O}_X}^\mathbf{L} \omega_{X/Y}^\bullet)
\ar@{=}[r] &
K \otimes_{\mathcal{O}_Y}^\mathbf{L} Rf_*\omega_{X/Y}^\bullet
\ar[rr]^-{\text{id}_K \otimes \text{Tr}_{f, \mathcal{O}_Y}} & & K
}
$$
where the equality on the lower right is
Derived Categories of Schemes, Lemma \href{https://stacks.math.columbia.edu/tag/08EU}{lemma cohomology base change (Tag 08EU)}.
If $g : Y' \to Y$ is a
morphism of quasi-compact and quasi-separated schemes
and $X' = Y' \times_Y X$, then by
Lemma \href{https://stacks.math.columbia.edu/tag/0AAB}{lemma proper flat base change (Tag 0AAB)} we have
$\omega_{X'/Y'}^\bullet = L(g')^*\omega_{X/Y}^\bullet$ where $g' : X' \to X$
is the projection and by Lemma \href{https://stacks.math.columbia.edu/tag/0B6J}{lemma trace map and base change (Tag 0B6J)}
the trace map
$$
\text{Tr}_{f', \mathcal{O}_{Y'}} :
Rf'_*\omega_{X'/Y'}^\bullet \to \mathcal{O}_{Y'}
$$
for $f' : X' \to Y'$ is the base change of $\text{Tr}_{f, \mathcal{O}_Y}$
via the base change isomorphism.
\end{remark}

\begin{lemma}
\label{lemma-properties-relative-dualizing}
Let $Y$ be a quasi-compact and quasi-separated scheme.
Let $f : X \to Y$ be a morphism of schemes which is
proper, flat, and of finite presentation with
relative dualizing complex $\omega_{X/Y}^\bullet$
(Remark \ref{remark-relative-dualizing-complex}).
Then
\begin{enumerate}
\item $\omega_{X/Y}^\bullet$ is a $Y$-perfect object of $D(\mathcal{O}_X)$,
\item $Rf_*\omega_{X/Y}^\bullet$ has vanishing cohomology sheaves
in positive degrees,
\item $\mathcal{O}_X \to
R\SheafHom_{\mathcal{O}_X}(\omega_{X/Y}^\bullet, \omega_{X/Y}^\bullet)$
is an isomorphism.
\end{enumerate}
\end{lemma}

\begin{proof}
In view of the fact that formation of $\omega_{X/Y}^\bullet$ commutes
with base change (see Remark \ref{remark-relative-dualizing-complex}),
we may and do assume that $Y$ is affine. For a perfect object $E$ of
$D(\mathcal{O}_X)$ we have
\begin{align*}
Rf_*(E \otimes_{\mathcal{O}_X}^\mathbf{L} \omega_{X/Y}^\bullet)
& =
Rf_*R\SheafHom_{\mathcal{O}_X}(E^\vee, \omega_{X/Y}^\bullet) \\
& =
R\SheafHom_{\mathcal{O}_Y}(Rf_*E^\vee, \mathcal{O}_Y) \\
& =
(Rf_*E^\vee)^\vee
\end{align*}
For the first equality, see
Cohomology, Lemma \href{https://stacks.math.columbia.edu/tag/08DQ}{lemma dual perfect complex (Tag 08DQ)}.
For the second equality, see Lemma \href{https://stacks.math.columbia.edu/tag/0A9Q}{lemma iso on RSheafHom (Tag 0A9Q)},
Remark \href{https://stacks.math.columbia.edu/tag/0GEV}{remark iso on RSheafHom (Tag 0GEV)}, and 
Derived Categories of Schemes, Lemma
\href{https://stacks.math.columbia.edu/tag/0B91}{lemma flat proper perfect direct image general (Tag 0B91)}.
The third equality is the definition of the dual. In particular
these references also show that the outcome is a perfect object
of $D(\mathcal{O}_Y)$. We conclude that $\omega_{X/Y}^\bullet$
is $Y$-perfect by More on Morphisms, Lemma
\href{https://stacks.math.columbia.edu/tag/0GET}{lemma characterize relatively perfect (Tag 0GET)}.
This proves (1).

\medskip\noindent
Let $M$ be an object of $D_\QCoh(\mathcal{O}_Y)$. Then
\begin{align*}
\Hom_Y(M, Rf_*\omega_{X/Y}^\bullet) & =
\Hom_X(Lf^*M, \omega_{X/Y}^\bullet) \\
& =
\Hom_Y(Rf_*Lf^*M, \mathcal{O}_Y) \\
& =
\Hom_Y(M \otimes_{\mathcal{O}_Y}^\mathbf{L} Rf_*\mathcal{O}_X, \mathcal{O}_Y)
\end{align*}
The first equality holds by Cohomology, Lemma
\href{https://stacks.math.columbia.edu/tag/079W}{lemma adjoint (Tag 079W)}.
The second equality by construction of $a$.
The third equality by Derived Categories of Schemes, Lemma
\href{https://stacks.math.columbia.edu/tag/08EU}{lemma cohomology base change (Tag 08EU)}.
Recall $Rf_*\mathcal{O}_X$ is perfect of tor amplitude in $[0, N]$
for some $N$, see
Derived Categories of Schemes, Lemma
\href{https://stacks.math.columbia.edu/tag/0B91}{lemma flat proper perfect direct image general (Tag 0B91)}.
Thus we can represent $Rf_*\mathcal{O}_X$ by a complex of
finite projective modules sitting in degrees $[0, N]$
(using More on Algebra, Lemma \href{https://stacks.math.columbia.edu/tag/0658}{lemma perfect (Tag 0658)}
and the fact that $Y$ is affine).
Hence if $M = \mathcal{O}_Y[-i]$ for some $i > 0$, then the last
group is zero. Since $Y$ is affine we conclude that
$H^i(Rf_*\omega_{X/Y}^\bullet) = 0$ for $i > 0$.
This proves (2).

\medskip\noindent
Let $E$ be a perfect object of $D_\QCoh(\mathcal{O}_X)$. Then
we have
\begin{align*}
\Hom_X(E, R\SheafHom_{\mathcal{O}_X}(\omega_{X/Y}^\bullet, \omega_{X/Y}^\bullet)
& =
\Hom_X(E \otimes_{\mathcal{O}_X}^\mathbf{L} \omega_{X/Y}^\bullet,
\omega_{X/Y}^\bullet) \\
& =
\Hom_Y(Rf_*(E \otimes_{\mathcal{O}_X}^\mathbf{L} \omega_{X/Y}^\bullet),
\mathcal{O}_Y) \\
& =
\Hom_Y(Rf_*(R\SheafHom_{\mathcal{O}_X}(E^\vee, \omega_{X/Y}^\bullet)),
\mathcal{O}_Y) \\
& =
\Hom_Y(R\SheafHom_{\mathcal{O}_Y}(Rf_*E^\vee, \mathcal{O}_Y),
\mathcal{O}_Y) \\
& =
R\Gamma(Y, Rf_*E^\vee) \\
& =
\Hom_X(E, \mathcal{O}_X)
\end{align*}
The first equality holds by Cohomology, Lemma
\href{https://stacks.math.columbia.edu/tag/08DJ}{lemma internal hom (Tag 08DJ)}.
The second equality is the definition of $\omega_{X/Y}^\bullet$.
The third equality comes from the construction of the dual perfect
complex $E^\vee$, see Cohomology, Lemma
\href{https://stacks.math.columbia.edu/tag/08DQ}{lemma dual perfect complex (Tag 08DQ)}.
The fourth equality follows from the equality
$Rf_*R\SheafHom_{\mathcal{O}_X}(E^\vee, \omega_{X/Y}^\bullet) =
R\SheafHom_{\mathcal{O}_Y}(Rf_*E^\vee, \mathcal{O}_Y)$
shown in the first paragraph of the proof.
The fifth equality holds by double duality for perfect complexes
(Cohomology, Lemma
\href{https://stacks.math.columbia.edu/tag/08DQ}{lemma dual perfect complex (Tag 08DQ)})
and the fact that $Rf_*E$ is perfect by
Derived Categories of Schemes, Lemma
\href{https://stacks.math.columbia.edu/tag/0B91}{lemma flat proper perfect direct image general (Tag 0B91)}.
The last equality is Leray for $f$.
This string of equalities essentially shows (3)
holds by the Yoneda lemma. Namely, the object
$R\SheafHom(\omega_{X/Y}^\bullet, \omega_{X/Y}^\bullet)$
is in $D_\QCoh(\mathcal{O}_X)$ by Derived Categories of Schemes, Lemma
\href{https://stacks.math.columbia.edu/tag/0A6H}{lemma quasi coherence internal hom (Tag 0A6H)}.
Taking $E = \mathcal{O}_X$ in the above we get a map
$\alpha : \mathcal{O}_X \to
R\SheafHom_{\mathcal{O}_X}(\omega_{X/Y}^\bullet, \omega_{X/Y}^\bullet)$
corresponding to
$\text{id}_{\mathcal{O}_X} \in \Hom_X(\mathcal{O}_X, \mathcal{O}_X)$.
Since all the isomorphisms above are functorial in $E$ we
see that the cone on $\alpha$ is an object $C$ of $D_\QCoh(\mathcal{O}_X)$
such that $\Hom(E, C) = 0$ for all perfect $E$.
Since the perfect objects generate
(Derived Categories of Schemes, Theorem
\href{https://stacks.math.columbia.edu/tag/09IS}{theorem bondal van den Bergh (Tag 09IS)})
we conclude that $\alpha$ is an isomorphism.
\end{proof}
\end{document}
